% Part: lambda-calculus
% Chapter: syntax
% Section: abbreviated-syntax

\documentclass[../../../include/open-logic-section]{subfiles}

\begin{document}

\olfileid[pt]{lam}{syn}{abb}
\olsection{Sintaxe abreviada}

Os termos definidos em \olref[trm]{defn:term} às vezes são incômodos de
escrever; por isso, é útil introduzir uma sintaxe mais concisa. Devemos,
naturalmente, ter o cuidado de garantir que os termos na notação concisa
também tenham leitura única. Uma versão amplamente usada, chamada
\emph{termos abreviados}, é a seguinte.

\begin{enumerate}
\item Quando os parênteses são omitidos, a aplicação é feita da esquerda
  para a direita. Por exemplo, se $M$, $N$, $P$ e~$Q$ são termos, então
  $MNPQ$ abrevia $(((MN)P)Q)$.
\item Novamente, quando os parênteses são omitidos, a abstração lambda recebe
  o maior escopo possível. Por exemplo, $\lambd[x][MNP]$ é lido como
  $(\lambd[x][MNP])$.
\item Um lambda pode ser usado para abstrair várias variáveis. Por exemplo,
  $\lambd[xyz][M]$ abrevia
  $\lambd[x][\lambd[y][\lambd[z][M]]]$.
\end{enumerate}

Por exemplo,
\[
\lambd[xy][xxyx \lambd[z][xz]]
\]
abrevia
\[
(\lambd[x][(\lambd[y][((((xx)y)x)(\lambd[z][(xz)]))])]).
\]

\begin{prob}
Expanda o termo abreviado $\lambd[g][(\lambd[x][g (x x)])
  \lambd[x][g (x x)]]$.
\end{prob}

\end{document}

